% Pre-Test 21 — Visual Strategies: Diagrams, number lines, pictures
\clearpage
\thispagestyle{empty}
{\sffamily

\begin{center}
{\fontsize{22}{26}\selectfont\bfseries\color{funPink} \faPalette\hspace{3mm}Draw It Out\hspace{3mm}\faPalette}

\vspace{2mm}
{\color{funGray}\hrulefill}
\end{center}

\vspace{3mm}

\begin{tcolorbox}[
    enhanced,
    colback=funPinkLight,
    colframe=funPink!50,
    arc=10pt,
    boxrule=1.5pt,
    left=6mm, right=6mm, top=5mm, bottom=5mm,
]
{\large\bfseries\color{funPinkDark} \faDrawPolygon\hspace{2mm}Pictures Make Problems Easier}

\vspace{2mm}

Some problems are hard to solve with numbers alone. Drawing a quick picture can make the answer obvious. Try these:

\vspace{3mm}

\textcolor{funPink}{\faRulerHorizontal}\hspace{2mm}\textbf{Number Lines} — Perfect for fractions, decimals, and integers. Plot the values and see which is bigger or smaller.

\vspace{2mm}

\textcolor{funPink}{\faShapes}\hspace{2mm}\textbf{Diagrams} — Geometry problems beg for a sketch. Draw the shape, label the sides, mark what you know.

\vspace{2mm}

\textcolor{funPink}{\faChartPie}\hspace{2mm}\textbf{Pie Charts \& Bar Models} — Great for ratio and percent problems. Divide a rectangle into parts to visualize the relationship.

\vspace{2mm}

\textcolor{funPink}{\faTable}\hspace{2mm}\textbf{Tables} — For rate and proportion problems, organize the information into rows and columns.
\end{tcolorbox}

\vspace{3mm}

\begin{tcolorbox}[
    enhanced,
    colback=funYellowLight,
    colframe=funYellow!60,
    arc=10pt,
    boxrule=1pt,
    left=6mm, right=6mm, top=4mm, bottom=4mm,
]
{\bfseries\color{funYellowDark} \faLightbulb\hspace{2mm}Try This}

\vspace{2mm}

On today's test, draw at least \textbf{3 pictures} on your scratch paper — even a quick sketch counts. See if it helps you solve faster or more accurately.
\end{tcolorbox}

\vspace{4mm}

\begin{center}
{\large\bfseries\color{funPinkDark} \faPaintBrush\hspace{2mm}A picture is worth a thousand calculations.\hspace{2mm}\faPaintBrush}
\end{center}

}
\newpage
